\documentclass[sigconf]{acmart}

% --- no ACM copyright block, conference banner or "ACM Reference Format" box (class report) ---
\setcopyright{none}
\settopmatter{printacmref=false}
\renewcommand\footnotetextcopyrightpermission[1]{}
\acmConference{}{}{}
\acmDOI{}
\acmISBN{}
\pagestyle{plain}

% Citations are numbered by default. For author-year citations, uncomment:
% \citestyle{acmauthoryear}

\begin{document}

\title{Disclosure Without Inference: LLM Agents Anchor on Disclosed Trade Data When Bargaining}

\author{Joey Poh}
\affiliation{%
  \institution{CS6101, National University of Singapore}
  \country{Singapore}}

\begin{abstract}
With the rapid improvement of LLM capabilities, an agentic economy, where autonomous AI agents act as independent actors that negotiate and settle transactions on behalf of humans, is becoming a reality. AGI Institutions has raised concerns over negotiations across a capability gap. % TODO: add a citation for AGI Institutions
Currently, asymmetric capabilities between negotiators are tolerated because the gap is bounded. However, in an agentic economy, asymmetric capabilities could potentially be much more extreme than the human case of a skilled negotiator versus an unskilled one, making it harder to find protocols that work for all parties. One design choice for such a negotiation is a floor mechanism. This is already seen in traditional financial markets: when the Financial Industry Regulatory Authority (FINRA) began publishing US corporate bond trades, investors negotiated better prices and their trading costs fell~\cite{goldstein2007transparency}. This paper investigates whether similar information would be effective for LLM agents.
\end{abstract}

\maketitle

% ======================================================================
\section{Introduction}

The experiment is a dyadic negotiation, where a client agent buys an off-the-run US Treasury note from a single dealer with a private cost. In one scenario, it also sees a summary of other dealers' reported purchases. We build a calibrated bargaining game with a Bayesian benchmark, under which the disclosure raises the client's share of the surplus by 11\% on our test scenarios. We then test a 35B open-weight model that crosses disclosure with belief elicitation.

Disclosure did not reliably improve the LLM's outcome (difference $-0.07$, 95\% CI $[-0.19, 0.04]$), well short of the Bayesian gain. The client did respond to the disclosure, but by bidding the numbers it was given: its counteroffers moved towards the disclosed typical cost and often sat on the bottom of the disclosed range, where most were rejected. Its stated beliefs did not drive its actions. In conclusion, an information floor cannot substitute for inference the agent does not perform.

% ======================================================================
\section{Background}

\paragraph{Bargaining with private values.}
When both sides of a trade hold private values, no bargaining procedure can guarantee that every worthwhile trade happens while giving both sides an incentive to be truthful~\cite{myerson1983efficient}. How the surplus is split then depends on the protocol and on what each side knows. Information design studies how a third party can change outcomes by choosing what to reveal~\cite{kamenica2011bayesian}. In corporate bonds, the introduction of trade reporting through TRACE lowered investors' trading costs~\cite{bessembinder2006market,edwards2007corporate}. Those studies measure market-wide averages, whereas we look inside a single negotiation.

\paragraph{LLM agents in bargaining.}
\citet{lewis2017deal} trained negotiation agents end to end. \citet{qian2026strategic} compared humans, LLMs and Bayesian agents in a multi-player trading game and found that similar surplus can hide different strategies: LLMs made concessionary offers that were usually accepted. Like them, we use induced values and a computed benchmark, but we hold the agent fixed and vary its information.

\paragraph{Bayesian theory of mind and anchoring.}
Inferring another agent's private state from its actions can be modelled as Bayesian inverse planning~\cite{baker2017rational,zhixuan2020online}. A client able to read a dealer's cost from its quote would need no disclosure. People, by contrast, often anchor on salient numbers~\cite{tversky1974judgment}, and LLMs show similar biases~\cite{jones2022capturing}. Our design distinguishes Bayesian use of a disclosure from anchoring on it.

% ======================================================================
\section{Bargaining Game with a Bayesian Benchmark}

In this game, a client buys \$10M of an off-the-run 10-year note from one dealer in a game limited to three moves. First, the dealer quotes a price $p$. The client either accepts $p$, or makes a counter-offer $q < p$. The dealer accepts $q$ if and only if $q \ge c$, its private cost.

The client has a private limit $r$, the most it will pay, and earns $r$ minus the price, or zero if there is no trade. Both sides see an indicative evaluated price $m$. Our primary outcome is the client's share of the surplus, zero if there is no trade, or its capture as defined as follows:
\begin{equation}
  \text{Capture} = \frac{r - \text{price}}{r - c}.
\end{equation}

\paragraph{Calibration.}
Each scenario draws the dealer's cost $c$ below $m$ and the client's limit $r$ above it, each at a random distance (LogNormal, median $1g$). Both sides know these distributions. The unit $g$, about \$3{,}400 on \$10mm, is one tenth of the note's daily price volatility, which puts it within the range of real dealer margins.

\paragraph{Disclosure.}
Ten other dealers' reported purchases are centred on the market-wide typical cost, from which this dealer deviates by a random amount $d \sim \mathcal{N}\!\left(0, (0.25g)^2\right)$. In the simulations where the band was disclosed, the client reads:

\begin{quote}
``Reported purchases by other dealers of similar off-the-run notes yesterday (10 trades) averaged $X$ below their evaluated prices. Applied to today's evaluated price, that implies a typical dealer cost of \textit{avg}. Individual dealers' costs vary around the typical level, so this dealer's own cost is within [\textit{low}, \textit{high}] about two times in three.''
\end{quote}

The band's centre is an unbiased but noisy reading of this dealer's cost, with standard deviation $s = 0.356g$, against a prior standard deviation of $0.604g$. The stated range, $\textit{avg} \pm s$, contains the true cost 68\% of the time, so the statement is exactly true.

\paragraph{Rational Client.}
As a benchmark, we compute the play of a Bayesian client. Suppose the client holds a belief $G(q)$, the probability that the dealer's cost is at or below $q$. Then its best counter $q^*$ is the counter that maximises $(r - q)\,G(q)$, and the highest quote it should accept is $\theta = r - \max_q\,(r - q)\,G(q)$. It accepts the quote $p$ if $p \le \theta$, and otherwise counters at $q^*$: a lower counter means a better price but a higher chance of rejection. Without the band, $G$ is the prior. With the band, $G$ is the Bayesian posterior.

\paragraph{The Dealer.}
The same fixed algorithm in every condition. It chooses the quote that maximises its expected profit over the prior on the client's limit, assuming an uninformed client. It rounds the quote to 0.001 points, and accepts any counter at or above its cost. It quotes $m - 0.425g$ in 88\% of scenarios, so its quote reveals little about its cost: the quote explains 18\% of the variation in the cost.

\paragraph{Value of disclosure.}
We play the rational client against this dealer at the true cost and limit. Over 1{,}000 scenarios, the band raises the client's surplus by 11.35\% of the total surplus. On our 50 test scenarios it raises capture from 0.621 to 0.732 ($+0.111$). This is the most a client that uses the band perfectly can gain. It is each side's best response to its stated beliefs, not an equilibrium.

% ======================================================================
\section{Experiments}

\subsection{Design}

\paragraph{Choosing the client.}
Before the experiment, we chose the client model with a short pilot on separate scenarios. The pilot checked basic competence, for example, that the model never counter-offers above its own limit. Our first choice, Qwen3.6-35B, failed this check. Ornith-1.5-35B passed and was used throughout. It runs at temperature 0.5 with reasoning off, so it must bargain without step-by-step working.

\paragraph{The Experiment.}
We varied two things: whether the client sees the band, and whether we ask it the belief question. The belief question asks for its probability that the dealer's cost lies below each of nine prices, once before it sees the quote and once after, so that we can compare its beliefs with a Bayesian client's.

% ----------------------------------------------------------------------
\subsection{LLM Agent Bargaining Results without Disclosure}

Without the band, the client almost always counters (in 98--99\% of negotiations), and far more aggressively than a rational client: on average $1.3g$ below $m$. About half of its counters are rejected, so it keeps 36--41\% of the surplus, against 62\% for the rational client.

Its counters sit on numbers it has been given. When not asked about the dealer's cost, 77 of its 98 counters land on one of two prices from the briefing: the typical dealer cost ($1g$ below $m$; 46 counters) or the low end of the stated range of costs ($1.9g$ below $m$; 31 counters; Figure~\ref{fig:counters}a). When asked to estimate the dealer's cost at nine prices, it also bids those prices: half of its counters sit at $1.25g$, $1.5g$ or $2g$ below $m$, which almost never happens otherwise.

Its stated beliefs do not drive its bids. They imply a counter about $0.47g$ below $m$, but they do not predict its actual counters ($\rho = 0.08$), and its counters do not rise with its own limit.

% ----------------------------------------------------------------------
\subsection{Evaluating if an Information Floor Helps}

The disclosure did not improve the LLM's outcomes. Without the cost question, the client kept 41\% of the surplus without the band and 33\% with it (Figure~\ref{fig:capture}). The difference, $-7$ points (95\% CI $-19$ to $+4$), is not reliably different from zero, so our main test is not supported. A rational client gains 11 points from the same band, which is outside this range: even at best, the LLM gained about a third of what the band offers. With the cost question, the band made no difference either ($-0.2$ points, 95\% CI $-10$ to $+10$).

\begin{figure}[t]
  \centering
  \includegraphics[width=0.85\columnwidth]{figures/fig2_results.pdf}
  \caption{Client capture without and with the band. Rational client: open markers; LLM client, without the cost question: filled markers. The band raises the rational client's capture by 0.11; the LLM's change is not reliably different from zero.}
  \label{fig:capture}
\end{figure}

% ----------------------------------------------------------------------
\subsection{How the LLM Agent Used the Band}

The client reacted to the band, but by bidding its numbers rather than reasoning with them. As predicted, its counteroffers moved towards the band's centre, by $0.16g$ without the belief question (95\% CI $[0.08, 0.24]$) and $0.26g$ with it ($[0.19, 0.33]$).

An exploratory look in Figure~\ref{fig:counters} shows this. Without the belief question, 26\% of counteroffers sat exactly on the bottom of the band's stated range, against 1\% without the band. The dealer's cost lies below that point only about 16\% of the time, so 77\% of these counteroffers were rejected: this is why the band left the client slightly worse off. The briefing's typical cost remained its most common bid, and its bids still ignored its own limit.

\begin{figure}[t]
  \centering
  \includegraphics[width=\columnwidth]{figures/fig3_counters.pdf}
  \caption{Where the client bids, without the belief question. (a)~Without the band (grey), counteroffers pile up on numbers in the briefing. (b)~With the band (blue), they move onto the band's own numbers, most often the bottom of its stated range (shaded).}
  \label{fig:counters}
\end{figure}

\paragraph{What the LLM Agent Believed.}
Eliciting the LLM client's beliefs revealed that it failed to update its belief on the dealer's cost and failed to act on that belief. The band changed the client's beliefs: its best guess of the dealer's cost moved $0.23g$ towards the band's typical cost, but its beliefs were no more accurate. Compared with the Bayesian client's, they were as far off with the band as without it (change in mean squared error $+0.012$, 95\% CI $[-0.004, 0.030]$). The client moved its guess towards the band, but only about a third as far as the Bayesian client would. The client also did not act on its own beliefs: its counters were unrelated to the counters its stated beliefs imply.

The client also did not weigh evidence by how informative it was. After seeing the dealer's quote, which reveals little about the cost, it concluded the cost was higher, yet it moved only partly on the band, which is far more informative. So, the LLM was not capable of using the information floor as a rational agent. It noticed the information, but neither weighed it correctly nor acted on what it believed.

\begin{figure}[t]
  \centering
  \includegraphics[width=\columnwidth]{figures/fig4_beliefs.pdf}
  \caption{Does the client's belief follow the band? Each marker is one negotiation with the belief question: the band's centre against the client's best guess of the dealer's cost (the median of its stated belief), before the quote, both relative to $m$ in units of $g$. Lines are least-squares fits. A Bayesian client tracks the band closely (black; slope 0.84). The LLM tracks it only weakly with the band (blue squares; slope 0.31) and not at all without it (grey circles; slope $-0.07$), so the blue slope comes from seeing the band. The guesses fall in rows because they almost always land on a price from the belief question or the briefing.}
  \label{fig:beliefs}
\end{figure}

% ----------------------------------------------------------------------
\subsection{Confirmation from Unseen Scenarios}

We repeated the experiment on 20 never-seen scenarios. Figure~\ref{fig:heldout} compares the two and confirms that the two main findings hold. First, the band again had no reliable effect on the client's share of the surplus. Disclosure is worth little on the new scenarios, so they cannot show whether the LLM falls short of the rational gain. Second, in Figure~\ref{fig:heldout}b, the client again moved its bids towards the band's centre, confirming that the LLM agent anchors on the numbers given.

We also asked whether the band changes what the client believes about the dealer's cost, as it would for a Bayesian client (Figure~\ref{fig:heldout}c). It barely does: the client's belief followed the band weakly in the main experiment (slope 0.31) and not at all on the new scenarios (0.00), against about 0.8 for a Bayesian client. The band changed where the client bid, but not what it believed.

\begin{figure}[t]
  \centering
  \includegraphics[width=\columnwidth]{figures/fig5_heldout.pdf}
  \caption{Main experiment (filled; 50 scenarios, two runs) against 20 held-out scenarios (open; one run), with 95\% CIs. Diamonds: rational or Bayesian client. (a)~The band's effect on capture. (b)~How far bids moved towards the band's centre. (c)~How closely the client's belief follows the band's centre (slope).}
  \label{fig:heldout}
\end{figure}

% ======================================================================
\section{Conclusion and Future Work}

The information floor (disclosing other dealers' trade data) did not reliably improve an LLM client's outcome. It delivered significantly less than a Bayesian client gains from the same information. While the client did not ignore the disclosure, it treated the disclosed numbers as prices to bid. These results suggest Ornith-1.5-35B does not have the inference needed to use it. Hence, an agent's fundamental capability remains the most important factor. Given that the LLM agent is anchored on the given numbers, how the information is published is perhaps of the same importance as whether to publish. Asking an agent about its beliefs is not neutral either. The prices in our question became anchors, and the agent's stated beliefs were not the beliefs it acted on.

This experiment was deliberately made simple. Future work should test whether the failure is specific to this setup: running the model with reasoning on, prompt engineering, and comparing model families and sizes would better investigate whether LLM agents are capable of using an information floor to close a capability gap.

% ======================================================================
\bibliographystyle{ACM-Reference-Format}
\bibliography{refs}

\end{document}
